\documentclass[11pt,a4paper]{article}
\usepackage[margin=28mm]{geometry}
\usepackage[T1]{fontenc}
\usepackage{lmodern}
\usepackage{amsmath,amssymb,amsthm,mathtools}
\usepackage{booktabs,array}
\usepackage{microtype}
\usepackage{graphicx}
\usepackage{xcolor}
\usepackage[colorlinks=true,linkcolor=blue!45!black,citecolor=blue!45!black,urlcolor=blue!45!black]{hyperref}
\usepackage{listings}
\setlength{\emergencystretch}{2em}
\setlength{\tabcolsep}{5pt}
\lstset{basicstyle=\ttfamily\footnotesize,breaklines=true,columns=fullflexible,
  keepspaces=true,frame=single,rulecolor=\color{black!20},xleftmargin=2mm,xrightmargin=2mm}
\newtheorem{theorem}{Theorem}[section]
\newtheorem{lemma}[theorem]{Lemma}
\newtheorem{proposition}[theorem]{Proposition}
\theoremstyle{remark}
\newtheorem{remark}[theorem]{Remark}
\newcommand{\F}{\mathbb F_{23}}
\newcommand{\Qt}{Q_{23}^{\mathrm t}}
\newcommand{\coeff}[2]{[#1]#2}
\newcommand{\card}[1]{\lvert#1\rvert}
\DeclareMathOperator{\ev}{ev}
\title{The independent domination number of the\texorpdfstring{\\}{ }23 by 23 toroidal queens graph:\texorpdfstring{\\}{ }a computer-assisted proof and Lean formalization}
\author{Weixin Yu\\Independent Researcher\\\href{mailto:iweixin1214@gmail.com}{\texttt{iweixin1214@gmail.com}}}
\date{November 8, 2026}
\hypersetup{pdftitle={The independent domination number of the 23 by 23 toroidal queens graph},
  pdfauthor={Weixin Yu},pdfsubject={A computer-assisted proof and Lean formalization}}

\begin{document}
\maketitle
\begin{abstract}
We give a reproducible computer-assisted proof that the independent domination
number of the $23\times23$ toroidal queens graph is $12$. A displayed set of
twelve queens gives the upper bound. For the lower bound, a finite-field
coefficient-extraction argument first excludes dominating sets with at most ten
queens. For eleven independent queens, translation to zero coordinate sums and
seven coefficients of a degree-$44$ covering polynomial impose restrictions on
the first three coefficients of four root polynomials. These restrictions reduce
the remaining obstruction to candidate lists containing $2774$ masks and a
finite check of $40\,620\,104$ triples. The entire implication from the chessboard
definitions to the exact minimum is implemented in Lean~4 with mathlib.
Three closed computations use \texttt{native\_decide}; the compiler and native
evaluation trust required by those steps is explicitly recorded. Two earlier
finite enumerations and checked LRAT refutations provide additional evidence,
but are not hypotheses of the Lean theorem. This manuscript makes no priority
claim for the numerical value and concerns this fixed board size only.
\end{abstract}

\noindent\textbf{Keywords:} toroidal queens graph; independent domination;
finite fields; polynomial method; exhaustive computation; Lean.

\section{Problem, result, and scope}
The vertices of $\Qt$ are the pairs $(r,c)\in\F^2$. Distinct vertices
$(r,c)$ and $(r',c')$ are adjacent when at least one of
\[
 r=r',\qquad c=c',\qquad r+c=r'+c',\qquad r-c=r'-c'
\]
holds in $\F$. A queen covers its own square and its adjacent vertices.
A set is \emph{independent} if no two of its distinct queens attack each other;
it is \emph{dominating} if every square is covered. The independent domination
number $i(\Qt)$ is the minimum cardinality of a set with both properties.

\begin{theorem}\label{thm:main}
The independent domination number of the $23\times23$ toroidal queens graph is
\[
 i(\Qt)=12.
\]
\end{theorem}

Toroidal queens graphs and their domination parameters were studied by
Burger, Mynhardt, and Cockayne~\cite{BMC2001}, and upper bounds were developed by
Mynhardt~\cite{Mynhardt2003}. OEIS sequence A279404 records the independent
domination parameter~\cite{OEIS}. The full text of~\cite{Mynhardt2003} was
checked: its main constructions concern orders divisible by $2$ or $3$, and
Table~2 on page~166 lists independent domination values through order $16$.
It does not determine the value at order $23$. The OEIS page obtained in this
check lists values through order $21$, with linked 2026 contributions for
orders $19$ and $20$. A bounded literature search on November~8, 2026 did not
locate an earlier proof of the exact value at order $23$. This is not an
exhaustive determination of priority, and no priority claim is made.

The proof has three parts. Section~\ref{sec:small} excludes ten or fewer queens
without assuming independence. Sections~\ref{sec:eleven}--\ref{sec:finite}
exclude eleven independent queens using algebra followed by finite computation.
Section~\ref{sec:upper} gives the twelve-queen construction.
Section~\ref{sec:lean} describes the formal statements and their precise trust
boundary. The result makes no assertion about an arbitrary board size $n$.

\section{A coefficient-extraction lemma}\label{sec:grid}
All coefficients and evaluations in this section lie in $\F$. For $n\geq0$,
put $\sigma_n=\sum_{t\in\F}t^n$, with the usual polynomial convention $t^0=1$.
Then
\begin{equation}\label{eq:powersums}
 \sigma_n=
 \begin{cases}
 -1,&n>0\text{ and }22\mid n,\\
 0,&\text{otherwise}.
 \end{cases}
\end{equation}
For $n=0$ this is $23=0$ in $\F$; for positive $n$, it follows by summing
powers in the cyclic multiplicative group of order $22$.

\begin{lemma}\label{lem:grid}
Let $P\in\F[x,y]$ have total degree at most $44$, and suppose that
$P(x,y)=0$ for every $(x,y)\in\F^2$. If $0\leq a,b\leq22$ and $a+b>22$, then
$\coeff{x^ay^b}{P}=0$.
\end{lemma}
\begin{proof}
Consider the linear functional
\[
 \Lambda_{a,b}(P)=\sum_{x\in\F}\sum_{y\in\F}
       x^{22-a}y^{22-b}P(x,y).
\]
On a monomial $x^iy^j$ its value is
$\sigma_{22-a+i}\sigma_{22-b+j}$. If this value is nonzero, both subscripts
are positive multiples of $22$. Thus $i=a+22u$ and $j=b+22v$ for integers
$u,v\geq0$. Since $i+j\leq44$ and $a+b>22$, we have $u=v=0$.
For $(i,j)=(a,b)$ the value is $(-1)^2=1$.
Consequently $\Lambda_{a,b}(P)=\coeff{x^ay^b}{P}$.
The assumed vanishing of $P$ makes the left side zero.
\end{proof}

For a finite set $Z\subseteq\F$, define its root polynomial by
\[
 L_Z(T)=\prod_{z\in Z}(T-z).
\]
For four sets $R,C,S,D\subseteq\F$, define the covering polynomial
\begin{equation}\label{eq:cover}
 P_{R,C,S,D}(x,y)=L_R(x)L_C(y)L_S(x+y)L_D(x-y).
\end{equation}
If $R,C,S,D$ contain, respectively, the row, column, sum-diagonal, and
difference-diagonal labels of a dominating set of queens, then this polynomial
vanishes on $\F^2$: a covering queen makes at least one of its four factors zero.

\section{Excluding ten queens or fewer}\label{sec:small}
\begin{proposition}\label{prop:ten}
Every dominating set in $\Qt$, whether independent or not, has at least eleven
queens.
\end{proposition}
\begin{proof}
Suppose that a dominating set has at most ten queens. Each of its four
projection sets has size at most ten. Enlarge each projection independently
to a subset $R,C,S,D$ of $\F$ of size exactly ten; this is possible since
$\card{\F}=23$. The polynomial in~\eqref{eq:cover} still vanishes on $\F^2$.
It has total degree at most $40$, and its homogeneous part of degree $40$ is
\[
 x^{10}y^{10}(x+y)^{10}(x-y)^{10}
   =x^{10}y^{10}(x^2-y^2)^{10}.
\]
All other terms have degree at most $39$. Hence
\[
 \coeff{x^{20}y^{20}}{P_{R,C,S,D}}
   =(-1)^5\binom{10}{5}=-252=1\quad\text{in }\F.
\]
Lemma~\ref{lem:grid}, with $(a,b)=(20,20)$, says that the same coefficient
must be zero. This contradiction proves the proposition.
\end{proof}

\section{Necessary conditions for eleven independent queens}\label{sec:eleven}
Suppose for contradiction that $A\subseteq\F^2$ is independent and dominating
and $\card A=11$. Write its elements as $(r_i,c_i)$, and put
\[
 R=\{r_i\},\quad C=\{c_i\},\quad S=\{r_i+c_i\},\quad D=\{r_i-c_i\}.
\]
Independence makes each of the four projection maps injective on $A$;
therefore every one of these sets has size $11$.

\subsection{Centering and root coefficients}
Translate all queens by
\[
 \left(-11^{-1}\sum_i r_i,-11^{-1}\sum_i c_i\right).
\]
Here $11^{-1}=21$ in $\F$. Translation preserves cardinality, independence,
and domination. We may therefore assume
\begin{equation}\label{eq:center}
 \sum_{r\in R}r=\sum_{c\in C}c=\sum_{s\in S}s=\sum_{d\in D}d=0.
\end{equation}
No queen is required to be at the origin after this translation.

For each of the four sets $Z$, write
\begin{equation}\label{eq:rootjet}
 L_Z(T)=T^{11}+u_ZT^9+v_ZT^8+E_Z(T),\qquad\deg E_Z\leq7.
\end{equation}
The missing $T^{10}$ term follows from~\eqref{eq:center}.
The coefficients can equivalently be read from
\begin{equation}\label{eq:reverse}
 B_Z(T)=\prod_{z\in Z}(1-zT):\qquad
 u_Z=\coeff{T^2}{B_Z},\quad v_Z=\coeff{T^3}{B_Z}.
\end{equation}
The elementary coefficient identity
\begin{equation}\label{eq:newton2}
 2\coeff{T^2}{B_Z}=\left(\sum_{z\in Z}z\right)^2-\sum_{z\in Z}z^2
\end{equation}
follows either by expanding the product or by induction when one root is
inserted. Under~\eqref{eq:center}, it gives $2u_Z=-\sum_{z\in Z}z^2$.

\subsection{Seven coefficients of the covering polynomial}
The polynomial $P=P_{R,C,S,D}$ has degree at most $44$ and vanishes on the
whole board. By Lemma~\ref{lem:grid}, its coefficients at
\begin{equation}\label{eq:positions}
 (22,20),(21,21),(20,22),\quad
 (22,19),(21,20),(20,21),(19,22)
\end{equation}
are all zero.

For $j\in\{2,3\}$, define the four homogeneous kernels
\begin{align*}
 K_{R,j}&=x^{11-j}y^{11}(x+y)^{11}(x-y)^{11},\\
 K_{C,j}&=x^{11}y^{11-j}(x+y)^{11}(x-y)^{11},\\
 K_{S,j}&=x^{11}y^{11}(x+y)^{11-j}(x-y)^{11},\\
 K_{D,j}&=x^{11}y^{11}(x+y)^{11}(x-y)^{11-j}.
\end{align*}
The degree-$42$ part of $P$ is $\sum_Zu_ZK_{Z,2}$ and its degree-$41$ part is
$\sum_Zv_ZK_{Z,3}$. Indeed, a single $u_Z$ term drops the degree by two,
a single $v_Z$ term drops it by three, a tail term drops it by at least four,
and two or more nonleading factors also drop it by at least four.
Thus no cross term has been omitted in degrees $42$ and $41$.

Straight binomial expansion gives the following coefficient matrices over $\F$.
The column order in both cases is $R,C,S,D$.
\begin{equation}\label{eq:matrix2}
 \begin{array}{c|rrrr}
  &R&C&S&D\\\hline
  x^{22}y^{20}&0&0&1&22\\
  x^{21}y^{21}&21&2&0&0\\
  x^{20}y^{22}&0&0&22&1
 \end{array}
\end{equation}
\begin{equation}\label{eq:matrix3}
 \begin{array}{c|rrrr}
  &R&C&S&D\\\hline
  x^{22}y^{19}&8&0&17&17\\
  x^{21}y^{20}&0&2&7&16\\
  x^{20}y^{21}&21&0&16&16\\
  x^{19}y^{22}&0&15&6&17
 \end{array}
\end{equation}
These are exact polynomial calculations, not tests on sampled queen sets.
The accompanying Lean files prove characteristic-free sparse expansions using
\texttt{ring}, then extract and reduce the specified coefficients.

The first two equations in~\eqref{eq:matrix2} give
\begin{equation}\label{eq:u-equal}
 u_R=u_C,\qquad u_S=u_D.
\end{equation}
The matrix in~\eqref{eq:matrix3} is invertible over $\F$. For example an inverse is
\begin{equation}\label{eq:inverse3}
 \begin{pmatrix}
 16&0&6&0\\
 0&17&0&7\\
 1&19&4&22\\
 1&4&4&1
 \end{pmatrix}.
\end{equation}
Consequently
\begin{equation}\label{eq:v-zero}
 v_R=v_C=v_S=v_D=0.
\end{equation}

There is one further geometric identity. Injectivity of the projections and
$(r+c)^2+(r-c)^2=2(r^2+c^2)$ imply
\[
 \sum_{s\in S}s^2+\sum_{d\in D}d^2
   =2\left(\sum_{r\in R}r^2+\sum_{c\in C}c^2\right).
\]
Using~\eqref{eq:newton2} and~\eqref{eq:u-equal}, and cancelling the nonzero
element $2$, gives
\begin{equation}\label{eq:u-double}
 u_S=u_D=2u_R=2u_C.
\end{equation}
In particular, all the conditions used below are necessary for an eleven-queen
independent dominating set. We make no symmetry assumption about that set.

\section{The finite obstruction}\label{sec:finite}
For $b\in\F$, let $\mathcal P_b$ be the family of sets $Z\subseteq\F$ satisfying
\begin{equation}\label{eq:profiles}
 \card Z=11,\quad \sum_{z\in Z}z=0,\quad
 \coeff{T^2}{B_Z}=b,\quad \coeff{T^3}{B_Z}=0.
\end{equation}
The preceding section gives $R,C\in\mathcal P_b$ and
$S,D\in\mathcal P_{2b}$ for $b=u_R$.

For any three sets of line labels, put
\begin{equation}\label{eq:delta}
 \Delta(R,C,S)=\{d\in\F:\exists r\in\F,
       r\notin R,\ r-d\notin C,\ 2r-d\notin S\}.
\end{equation}
Equivalently these are the difference labels of squares left uncovered by the
row, column, and sum-diagonal families. If the four line families cover the
board, then
\begin{equation}\label{eq:subset}
 \Delta(R,C,S)\subseteq D.
\end{equation}
Indeed, for a witness $r$ in~\eqref{eq:delta}, the square $(r,r-d)$ has no
covering row, column, or sum diagonal, so must be covered by difference label $d$.

\begin{lemma}[finite computation]\label{lem:finite}
For all $b\in\F$ and all $R,C\in\mathcal P_b$, $S\in\mathcal P_{2b}$,
\[
 \card{\Delta(R,C,S)}\geq12.
\]
\end{lemma}

Here the finite computation is part of the proof. Its definition, exhaustive
scope, and connection to the stated sets are described next.

\subsection{Candidate lists and certified completeness}
Encode $Z\subseteq\F$ as the natural number
$m(Z)=\sum_{z\in Z}2^{\operatorname{val}(z)}$, where residues are represented
by $0,\ldots,22$. The file \texttt{ProfileData.lean} contains $23$ explicit
candidate lists $\mathcal L_b$.
To prove that every profile is in the appropriate list, an executable Lean
function scans \emph{every} mask $0\leq m<2^{23}=8\,388\,608$.

For a partial product $\prod(1-zT)$, keep a state $(k,a,b,c)$ consisting of
the number of roots and coefficients of $T,T^2,T^3$. Inserting a root $z$
updates this state by
\begin{equation}\label{eq:recurrence}
 (k,a,b,c)\longmapsto(k+1,a-z,b-za,c-zb).
\end{equation}
The initial state is $(0,0,0,0)$, and a root is inserted precisely when the
corresponding bit of $m$ is set. For every mask with $k=11$, $a=0$, and $c=0$,
the checker requires $m\in\mathcal L_b$.
The scan is evaluated by \texttt{native\_decide}. An ordinary Lean induction
proves that this recurrence equals the coefficients of~\eqref{eq:reverse},
and another proves that a successful range scan checks each integer in its range.
Together these establish
\begin{equation}\label{eq:list-complete}
 Z\in\mathcal P_b\ \Longrightarrow\ m(Z)\in\mathcal L_b.
\end{equation}
The proof needs this inclusion only; it does not assume that external list
generation is correct or require the converse inclusion as a hypothesis.

There are $2774$ entries in the lists. Their sizes and the resulting triple
counts are shown in Table~\ref{tab:counts}. The coefficient parameter $b$ is
not the second power sum: when the sum is zero,
$b=-\frac12\sum z^2=11\sum z^2$ in $\F$.

\begin{table}[ht]
\centering
\begin{tabular}{lrrr}
\toprule
Coefficient parameter $b$ & Number of groups & Entries per group & Triples in all groups\\
\midrule
$0$ & 1 & 134 & $2\,406\,104$\\
Nonzero quadratic residue & 11 & 115 & $16\,729\,625$\\
Quadratic nonresidue & 11 & 125 & $21\,484\,375$\\
\midrule
Total & 23 & $2774$ entries & $40\,620\,104$\\
\bottomrule
\end{tabular}
\caption{Candidate list sizes and exhaustive coverage checks. Since $2$ is a
square modulo $23$, $b$ and $2b$ have the same group size.}
\label{tab:counts}
\end{table}

\subsection{The executable coverage check and its semantics}
Masks in this computation are $23$-bit vectors. If $M$ is a mask,
$\operatorname{rot}_t(M)$ denotes cyclic left rotation by $t$ places;
it represents the set obtained by adding $t$ modulo $23$.
Let $H(S)$ be the mask whose bit at $r$ is the bit of $S$ at $2r$.
It represents $S/2$, since $2^{-1}=12$ in $\F$.
For each $d\in\{0,\ldots,22\}$, form the mask
\begin{equation}\label{eq:bitcheck}
 U_d=\overline R\ \mathbin{\&}\ 
       \overline{\operatorname{rot}_d(C)}\ \mathbin{\&}\ 
       \overline{\operatorname{rot}_{12d}(H(S))}.
\end{equation}
Complements are taken within $23$ bits.
Its bit at $r$ is set exactly when
$r\notin R$, $r-d\notin C$, and $2r-d\notin S$.
Therefore $d\in\Delta(R,C,S)$ if and only if $U_d\ne0$.

The executable checker loops over all $23$ values of $b$, every column mask
in $\mathcal L_b$, every sum mask in $\mathcal L_{2b}$, and every row mask in
$\mathcal L_b$. For each triple it counts the $23$ indices with $U_d\ne0$
and requires the count to be at least $12$. A closed
\texttt{native\_decide} theorem certifies that every check returns true.
Rotations and allowed-row masks are cached outside the inner row loop.
All groups are checked; the Lean proof uses no scaling reduction.

The file \texttt{CoverSemantics.lean} proves that this executable count is
exactly the cardinality of~\eqref{eq:delta}. In particular, the proof relates
bit-vector rotations, residue subtraction, set encoding, and finite-set
cardinality. Applying this bridge to the successful exhaustive check and
using~\eqref{eq:list-complete} proves Lemma~\ref{lem:finite}.

\subsection{The contradiction}
For the hypothesized eleven-queen set, $\card D=11$ and~\eqref{eq:subset} give
$\card{\Delta(R,C,S)}\leq11$. Lemma~\ref{lem:finite} gives the opposite bound
$\card{\Delta(R,C,S)}\geq12$. This excludes eleven independent dominating
queens. With Proposition~\ref{prop:ten}, every independent dominating set has
at least twelve queens.

\section{A twelve-queen construction}\label{sec:upper}
All coordinates below are zero indexed and interpreted modulo $23$.
Take the twelve squares
\begin{equation}\label{eq:witness}
\begin{split}
 &(2,13),(3,9),(6,3),(7,22),(9,2),(11,5),\\
 &(12,8),(15,18),(16,14),(18,7),(20,4),(21,0).
\end{split}
\end{equation}
There are twelve distinct rows, columns, sum labels, and difference labels,
so the queens are independent. Direct evaluation on the $529$ squares shows
that they dominate the board. Both assertions, and cardinality twelve, are
proved by ordinary \texttt{decide} in \texttt{Upper.lean}; this upper-bound
proof does not use \texttt{native\_decide}. Combined with the lower bound,
it proves Theorem~\ref{thm:main}.

\section{Lean implementation and trust boundary}\label{sec:lean}
\subsection{Formal statements and dependency structure}
The implementation uses Lean $4.34.0$~\cite{Lean} and mathlib~\cite{mathlib} source revision
\begin{center}
\texttt{5ed2965256430c3649e86755f9576b54eca72435}.
\end{center}
The proof models the field as \texttt{ZMod 23} and squares as pairs of field
elements. Its definitions of independence and domination quantify over the
actual squares and queens, rather than assuming any polynomial equations or
profile restrictions. The main exact-minimum statement is
\begin{lstlisting}
theorem minimum_is_twelve :
    (exists s : Finset Square,
      s.card = 12 /\ Independent s /\ Dominating s) /\
    (forall s : Finset Square,
      Independent s -> Dominating s -> 12 <= s.card)
\end{lstlisting}
Here the ASCII display gives the logical content of the Unicode source.
The source also defines the minimum by the infimum of attainable cardinalities
and proves the equality
\begin{lstlisting}
theorem independentDominationNumber_eq_twelve :
    independentDominationNumber = 12
\end{lstlisting}
Both the existence and universal lower bound are unconditional within the
declared Lean/native computation trust boundary.

The modules have the following roles.
\begin{table}[ht]
\centering\small
\begin{tabular}{>{\raggedright\arraybackslash}p{.33\linewidth}>{\raggedright\arraybackslash}p{.59\linewidth}}
\toprule
Module(s) & Content\\
\midrule
\texttt{Basic} & Chessboard definitions, translations, centering, projection injectivity.\\
\texttt{RootCoefficients} & Root products, reversal, coefficient identities and tails.\\
\texttt{Grid}, \texttt{Lift} & Grid coefficient extraction, substitution and degree bounds.\\
\texttt{Kernels}, \texttt{Expansion}, \texttt{Equations} & Exact expansions, error-degree bounds, seven coefficients and linear elimination.\\
\texttt{Geometry}, \texttt{LowerTen} & Projection identities, coefficient restrictions, and exclusion of at most ten queens.\\
\texttt{ProfileData}, \texttt{ProfileEnumeration}, \texttt{Masks}, \texttt{ProfileBridge} & Explicit lists, exhaustive completeness, mask semantics and coefficient recurrence.\\
\texttt{CoverCheck}, \texttt{CoverSemantics} & Exhaustive triple check and its interpretation as $\card\Delta$.\\
\texttt{Upper}, \texttt{Main} & Twelve-queen witness and the complete minimum theorem.\\
\bottomrule
\end{tabular}
\caption{Structure of the accompanying Lean source. Module names are relative
to the namespace \texttt{ToroidalQueens23}.}
\end{table}

\subsection{Native computation is explicitly trusted}
The final build passed with no \texttt{sorry}, \texttt{admit}, or manually
declared axiom in the proof source. The algebraic identities, geometric
bridges, recurrences, and range-scan soundness are ordinary Lean proofs.
Exactly three closed calculations use \texttt{native\_decide}:
\begin{enumerate}
\item \texttt{powerSum\_table}: the power sums for $0\leq n<89$;
\item \texttt{all\_profile\_masks\_checked}: completeness over all $2^{23}$ masks;
\item \texttt{all\_profile\_covers\_checked}: all $23$ candidate groups and their
$40\,620\,104$ triples.
\end{enumerate}
The axiom report for the final equality includes \texttt{propext},
\texttt{Classical.choice}, and \texttt{Quot.sound}, together with the three
generated native computation axioms named in Appendix~\ref{app:audit}.
There is no \texttt{sorryAx} in that report.

This is a complete Lean-compiled formalization with a native computation
trust component. It is not claimed to be a proof reduced entirely by the
Lean kernel using only the three standard axioms. Correctness of the three
native steps additionally relies on the relevant compiler/runtime and native
evaluation mechanism. The finite computations are closed Lean programs with
proved semantic bridges; no Python result or SAT solver status is imported
as an assumed theorem.

\subsection{Supplementary checks and reproducibility}
Before the complete formalization, two different implementations checked the
finite obstruction independently. One enumerated all
$\binom{23}{11}=1\,352\,078$ eleven-element subsets and accumulated uncovered
difference labels directly. The other used fixed-weight integer masks and
vectorized integer rotations. Both checked all $40\,620\,104$ triples; the
second reported a minimum of $13$ required difference labels, and $15$ in the
zero-parameter group. These stronger observed minima are supplementary
results, not the bound certified or needed by the Lean theorem.

Three additional profile SAT instances, for second-moment representatives
$0,1,22$, yielded LRAT refutations. A checker separate from the solver
verified all three to the empty clause, totaling $1\,867\,319$ clause additions
and $139\,363\,759$ positive hint references. These historical certificates and
their logs are retained. The Lean proof instead checks every coefficient
group directly and does not formalize or depend on those LRAT certificates.
A known five-queen construction on the $11\times11$ torus was also used as
a positive control for the earlier programs.

The source distribution includes the proof modules, exact candidate data,
toolchain pin, a portable Lake configuration, build/axiom audit, and checksums.
The manuscript source uses standard LaTeX packages. Rebuilding the formal
proof with a fresh mathlib environment requires installing the pinned Lean
toolchain, fetching dependencies, and running \texttt{lake build}. A successful
build re-evaluates the closed native checks when their modules are rebuilt;
on the Windows development environment, the coverage computation took about
$12$ minutes. Timings depend on hardware and are not proof hypotheses.
The checksums identify the exact tested source, and the build log distinguishes
newly compiled modules from cached ones. Historical failed development logs
are preserved separately from the final successful audit.

\section{Limitations and further work}
The result is a fixed finite instance of independent domination on a torus.
The covering-polynomial restrictions are necessary conditions; the proof does
not classify all optimal configurations, prove a general formula in $n$, or
establish a geometric symmetry of hypothetical critical configurations.
The current finite obstruction is computational. Replacing it by a short
structural proof, or replacing native evaluation by compact kernel-checkable
certificates, would strengthen the proof's explanatory or trust properties.

This manuscript is a preprint draft. The numerical value's literature priority
has not been established, and no external subject-expert review is claimed.
The proof and files were prepared with assistance from OpenAI Codex for
mathematical development, programming, formalization, and writing. Such
assistance or agreement is not evidence in place of the proofs and executable
checks specified above.

\appendix
\section{Exact native axiom audit}\label{app:audit}
With the common prefix \texttt{ToroidalQueens23.} omitted, the generated native
names appearing in the final equality's axiom report are
\begin{lstlisting}
all_profile_covers_checked._native.native_decide.ax_1_1
all_profile_masks_checked._native.native_decide.ax_1_1
powerSum_table._native.native_decide.ax_1_1
\end{lstlisting}
The remaining names are \texttt{propext}, \texttt{Classical.choice}, and
\texttt{Quot.sound}. The upper bound alone has only those three standard names.
The source bundle includes the literal report rather than treating this prose
as its proof certificate.

\section{Coefficient formula for reproducing the kernel tables}
For nonnegative integers $a,b,c,d$, the coefficient of $x^iy^j$ in
$x^ay^b(x+y)^c(x-y)^d$ is
\[
 \sum_{\substack{0\leq k\leq c,\ 0\leq\ell\leq d\\
                 a+c+d-k-\ell=i,\ b+k+\ell=j}}
     \binom ck\binom d\ell(-1)^\ell.
\]
Apply this formula with the exponent tuples
\[
 (11-j,11,11,11),\ (11,11-j,11,11),\
 (11,11,11-j,11),\ (11,11,11,11-j),
\]
for $j=2,3$, and reduce the resulting integers modulo $23$.
The Lean source proves the full integer-coefficient expansions by ring
normalization before extracting the table entries, so the generator that
emits these expressions is not a trusted oracle.

\begin{thebibliography}{9}
\raggedright
\bibitem{BMC2001}
A.~P. Burger, C.~M. Mynhardt, and E.~J. Cockayne,
\emph{Queens graphs for chessboards on the torus},
Australasian Journal of Combinatorics \textbf{24} (2001), 231--246.
\url{https://ajc.maths.uq.edu.au/pdf/24/ocr-ajc-v24-p231.pdf}.

\bibitem{Mynhardt2003}
C.~M. Mynhardt,
\emph{Upper bounds for the domination numbers of toroidal queens graphs},
Discussiones Mathematicae Graph Theory \textbf{23} (2003), 163--175.
\href{https://doi.org/10.7151/dmgt.1193}{doi:10.7151/dmgt.1193}.

\bibitem{OEIS}
The OEIS Foundation Inc.,
\emph{A279404: Independent domination number for queens' graph on an
$n\times n$ toroidal board}, accessed November 8, 2026.
\url{https://oeis.org/A279404}.

\bibitem{Lean}
The Lean development team,
\emph{Lean 4}, version 4.34.0.
\url{https://github.com/leanprover/lean4}.

\bibitem{mathlib}
The mathlib community,
\emph{mathlib4}, source revision
\texttt{5ed2965256430c3649e86755f9576b54eca72435}.
\url{https://github.com/leanprover-community/mathlib4}.
\end{thebibliography}
\end{document}
